\documentclass{article}
\usepackage{amsmath}
\usepackage[ruled,vlined,linesnumbered]{algorithm2e}

\begin{document}

\section*{La logica dietro la struttura heap}
L'algoritmo usa la struttura dati \textbf{\textit{heap}} per poter estrarre rapidamente
il valore massimo ($\Theta(\log n)$) sfruttando la logica del \textit{max heap},
ovvero la \textit{chiave} di un qualsiasi nodo $v$ sarà sempre maggiore della chiave dei suoi figli. \\
Un heap è un \textbf{albero binario} completo, quindi ogni nodo avrà due figli, fino al 
penultimo livello, ed è una struttura dati rappresentata tramite \textbf{vettore di posizione},
quindi preso un qualsiasi nodo $v$ per trovare i suoi figli seguiremo la solita formula 
$P[d \cdot v + i]$, che possiamo specificare come $P[2 \cdot v]$ per il figlio destro e 
$P[2 \cdot v + 1]$ per il sinistro. 

\section*{L'algoritmo}

% HEAPSORT %
\vspace{0.3cm}
\begin{algorithm}[H]
\SetKwFunction{FBuildMaxHeap}{buildMaxHeap}
\SetKwFunction{FHeapify}{heapify}
\SetKwProg{Fn}{heapsort}{}{}%

\Fn{$(array\ A)$}{
    \FBuildMaxHeap{$A$}\;

    \BlankLine

    \tcp{estrazione dei valori dall'heap}
    \For{$i = n$ \KwTo $1$}{
        \textbf{swap} $A[i], A[1]$\;
        $n \gets n - 1$\;
        \FHeapify{$A, 1$}\;
    }
}
\caption{heapsort}
\end{algorithm}
\vspace{0.3cm}

L'algoritmo presenta due funzioni ausiliarie importantissime che andremo ad analizzare 
nel dettaglio

% BUILD MAX HEAP %
\vspace{0.3cm}
\begin{algorithm}[H]
\SetKwFunction{FHeapify}{heapify}
\SetKwProg{Fn}{buildMaxHeap}{}{}%

\Fn{$(array\ A)$}{
    $n \gets A.len()$\;

    \BlankLine

    \tcp{costruzione dell'heap solo sui nodi padri per}
    \tcp{salvarci del tempo}
    \For{$i = \lfloor n/2 \rfloor$ \KwTo $1$}{
        \FHeapify{$A, i$}\;
    }
}
\caption{buildMaxHeap}
\end{algorithm}
\vspace{0.3cm}

$buildMaxHeap$ si occupa della costruzione della struttura dati dal semplice array iniziale. 
Notare come non siamo iterando su ogni nodo ma soltanto sulla prima metà perchè la funzione succesiva,
$heapify$, considera ogni figlio succesivo.

% HEAPIFY % 
\vspace{0.3cm}
\begin{algorithm}[H]
\SetKwFunction{FHeapify}{heapify}
\SetKwProg{Fn}{heapify}{}{}%

\Fn{$(array\ A, int\ father)$}{
    $newFather \gets father$\;
    $left \gets father \times 2$\;
    $right \gets father \times 2 + 1$\;

    \BlankLine
    \BlankLine

    \tcp{se l'indice left ha "sforato" la lunghezza dell'array}
    \tcp{e se è maggiore del padre}
    \If{$left \le A.len() \textbf{ and } A[left] > A[father]$}{
        $newFather \gets left$\;
    }

    \BlankLine

    \tcp{se l'indice right ha "sforato" la lunghezza dell'array}
    \tcp{e se è maggiore del padre}
    \If{$right \le A.len() \textbf{ and } A[right] > A[father]$}{
        $newFather \gets right$\;
    }

    \BlankLine

    \tcp{se il padre non è cambiato}
    \If{$newFather \neq father$}{
        \textbf{swap} $A[fater], A[newFather]$\;
        \FHeapify{$A, newFather$}\;
    }
}
\caption{heapify}
\end{algorithm}
\vspace{0.3cm}

Come già detto prima l'algoritmo sfrutta la formula $P[2 \cdot v]$ e $P[2 \cdot v + 1]$ per trovare
subito i figli del padre su cui stiamo attualmente lavorando. L'array su cui stiamo lavorando
non è ovviamente ordinato quindi i due $if$ servono proprio per gestire l'ordine tra i padri ed 
i figli.

\section*{Costo computazionale}
Il costo verrà ovviamente dettato dalla somma delle due funzioni, quindi:

\[ T(n) = \text{BuildMaxHeap} + \text{Heapify in Heapsort} \]

Come detto prima $buildMaxHeap$ itera usando un for a "metà", ma nel caso pegiore quell'altra
metà sarà analizzata da Heapify, portando il costo a $O(n)$. Infine notiamo come la funzione 
originale $heapsort$ itera usando un for da $n$ a $1$, ma non esegue esattamente $n$ cicli poichè
ad ogni iterazione viene si decrementato $i = i - 1$, ma succede lo stesso con $n = n - 1$, portando
quindi ad avere $\log n$ cicli, in cui al suo interno avremo un $heapify$ con costo $O(n)$.

\[ T(n) = \text{BuildMaxHeap} + \text{Heapify in Heapsort} \]
\[ \downarrow \]
\[ T(n) = O(n) + O(n \cdot \log n) = O(n \cdot \log n)\]

\end{document}